% BHU PYQ -- Manually transcribed & error-corrected
\documentclass[11pt,a4paper]{article}

\usepackage[a4paper,top=2.5cm,bottom=2.5cm,left=2.8cm,right=2.8cm,headheight=14pt]{geometry}
\usepackage{amsmath,amssymb}
\usepackage[T1]{fontenc}
\usepackage[utf8]{inputenc}
\usepackage{lmodern}
\usepackage{microtype}
\usepackage{fancyhdr}
\usepackage{enumitem}
\usepackage{hyperref}
\usepackage{lastpage}
\usepackage{parskip}

\hypersetup{colorlinks=true,urlcolor=black,linkcolor=black,
  pdftitle={BPT-401: Electronics and Modern Physics -- BHU PYQ},pdfauthor={Banaras Hindu University}}

\pagestyle{fancy}\fancyhf{}
\renewcommand{\headrulewidth}{0.4pt}\renewcommand{\footrulewidth}{0.4pt}
\fancyhead[L]{\small\textit{Banaras Hindu University}}
\fancyhead[R]{\small\textit{BPT-401: Electronics and Modern Physics}}
\fancyfoot[C]{\small Page \thepage\ of \pageref{LastPage}}

\newlist{parts}{enumerate}{2}
\setlist[parts,1]{label=(\alph*),leftmargin=*,itemsep=5pt,topsep=3pt}
\setlist[parts,2]{label=(\roman*),leftmargin=*,itemsep=3pt,topsep=2pt}
\newcommand{\pts}[1]{\hfill\textit{[#1]}}

\begin{document}

\begin{center}
  {\large\bfseries BANARAS HINDU UNIVERSITY}\\[2pt]
  {\normalsize B.Sc. (Hons.) Semester IV Examination, 2022-23}\\[6pt]
  \rule{0.8\linewidth}{0.5pt}\\[6pt]
  {\large\bfseries Paper No. BPT-401: Electronics and Modern Physics}\\[4pt]
  {\normalsize Time: 3 Hours\hfill Maximum Marks: 70}
\end{center}
\rule{\linewidth}{0.4pt}
\smallskip
\noindent\textit{Instructions: Answer \textbf{five} questions including Question~1 which is compulsory. Symbols have their usual meanings.}

\medskip\noindent\textbf{Question 1.} \textit{(Compulsory --- Answer any five parts, 4 marks each.)}
\begin{parts}
  \item What are the differences between metals, insulators, and semiconductors? Describe them on the basis of the band theory of solids. \pts{4}
  \item Define the phase velocity and group velocity of matter waves. Obtain the relation between them. \pts{4}
  \item If a dielectric is introduced between the plates of a parallel plate capacitor, show how the induced charge varies as a function of the dielectric constant. Show that the dielectric constant approaches infinity for a metal. \pts{4}
  \item Differentiate between the electronic, ionic, and orientational polarizabilities of a dielectric material. \pts{4}
  \item For an intrinsic semiconductor, the effective masses of hole and electron are $0.36\,m_0$ and $0.54\,m_0$ respectively ($m_0$ being the rest mass of an electron). Find the position of the Fermi level in this semiconductor at $T = 300\,\text{K}$. \pts{4}
  \item Light of wavelength $400\,\text{nm}$ is incident on a metal surface with a work function of $2.5\,\text{eV}$. Calculate the maximum kinetic energy of the emitted photoelectrons. \pts{4}
  \item How is the current amplification factor $\alpha$ different from the current gain $\beta$? Determine the current gain $\beta$ and emitter current $I_E$ for a transistor operating in CE configuration, where $I_B = 50\,\mu\text{A}$ and $I_C = 3.65\,\text{mA}$. \pts{4}
\end{parts}

\medskip\rule{\linewidth}{0.2pt}

\medskip\noindent\textbf{Question 2.}
\begin{parts}
  \item Derive the expression for the local molecular field inside a spherical cavity of a dielectric material and establish the Clausius-Mossotti equation. What are the limitations of this equation? \pts{7}
  \item Explain the three electric vectors: electric field intensity $\vec{E}$, dielectric polarization $\vec{P}$, and electric displacement $\vec{D}$. Establish the relation among them. \pts{7}
\end{parts}
\hfill \textbf{OR}
\begin{parts}
  \item Derive the Langevin-Debye formula for the dielectric susceptibility of a polar medium. \pts{7}
  \item Describe the formation of the depletion layer in an unbiased p-n junction. Sketch the barrier potential and show that the barrier potential is given by:
  \[
  V_0 = \frac{kT}{q} \ln \left( \frac{N_a N_d}{n_i^2} \right)
  \]
  where the symbols have their usual meanings. \pts{7}
\end{parts}

\medskip\rule{\linewidth}{0.2pt}

\medskip\noindent\textbf{Question 3.}
\begin{parts}
  \item What is meant by the peak inverse voltage (PIV) of a rectifier? Prove that the regulation for both half-wave and full-wave rectifiers is given by the ratio of forward diode resistance to the load resistance. \pts{5.5}
  \item Explain the difference between the Zeeman effect and the Compton effect. Give a brief account of Einstein's explanation of the photoelectric effect on the basis of quantum theory. \pts{5}
  \item Show that the de Broglie wavelength ($\lambda$) for a material particle of rest mass $m_0$ and charge $q$ accelerated from rest through a potential difference of $V$ volts relativistically is given by:
  \[
  \lambda = \frac{h}{\sqrt{2 m_0 q V \left(1 + \frac{q V}{2 m_0 c^2}\right)}}
  \]
  where $c$ is the speed of light. \pts{3.5}
\end{parts}
\hfill \textbf{OR}
\begin{parts}
  \item "Zener breakdown voltage has a negative temperature coefficient whereas the avalanche breakdown voltage has a positive temperature coefficient." Explain this statement. How does a Zener diode work as a voltage regulator? \pts{7}
  \item Derive the expression for the change in wavelength of a photon due to Compton scattering with a stationary electron. Explain the significance of the Compton wavelength. \pts{7}
\end{parts}

\medskip\rule{\linewidth}{0.2pt}

\medskip\noindent\textbf{Question 4.}
\begin{parts}
  \item State Bragg's law of X-ray diffraction in crystals. Describe the Bragg X-ray spectrometer method for determining the wavelength of X-rays. \pts{5}
  \item What is the meaning of biasing a transistor? Discuss the different methods used for transistor biasing and provide a neat circuit diagram for a self-biased transistor. \pts{6}
  \item Explain the terms "polar vector" and "axial vector" with one example of each. \pts{3}
\end{parts}
\hfill \textbf{OR}
\begin{parts}
  \item A silicon n-p-n transistor having $\beta = 100$, $I_{CO} = 22\,\text{nA}$ is operated in CE configuration. Assuming $V_{BE} = 0.7\,\text{V}$ and collector supply voltage $V_{CC} = 12\,\text{V}$ with collector load resistance $R_C = 2\,\text{k}\Omega$, determine the transistor currents and the region of operation. What happens if $R_C$ is indefinitely increased? \pts{10}
  \item What is a load line? Explain using a suitable diagram how the DC load line can be constructed for a common emitter transistor amplifier. \pts{4}
\end{parts}

\medskip\rule{\linewidth}{0.2pt}

\medskip\noindent\textbf{Question 5.}
\begin{parts}
  \item Describe the operation of an R-C coupled common-emitter transistor amplifier and explain its frequency response curve. \pts{7}
  \item What is the Early effect in transistors? Explain why the transistor action is lost if the base width is reduced to zero. \pts{7}
\end{parts}
\hfill \textbf{OR}
\begin{parts}
  \item Show that the energy lost by a photon in Compton scattering with a free stationary electron can be written as:
  \[
  \Delta E = \frac{(h\nu)^2 (1 - \cos\theta)}{m_0 c^2 + h\nu(1 - \cos\theta)}
  \]
  where $m_0$ is the rest mass of the electron, $h\nu$ is the energy of the incident photon, and $\theta$ is the scattering angle of the photon. \pts{7}
  \item Calculate the maximum percentage of energy that can be lost by a $1\,\text{MeV}$ photon in a collision with a stationary electron (rest mass energy of electron is $0.51\,\text{MeV}$). \pts{7}
\end{parts}

\vfill
\rule{\linewidth}{0.4pt}
\begin{center}\small
  \href{https://bhu-exam.vercel.app/}{Click here to visit our website}\\[4pt]
  \href{https://me-aryan.vercel.app/}{Click here to visit our contributor}\\[4pt]
  \href{https://quashan-me.vercel.app/}{Click here to visit our contributor}
\end{center}

\end{document}
